\section{Introduction}

A random projection maps $d$-dimensional data to $k$ dimensions with a random matrix chosen without looking at the data. The Johnson--Lindenstrauss lemma says that $k$ on the order of $\log n$ preserves the distances among $n$ points up to a small distortion \cite{johnson1984,dasgupta2003}. Achlioptas showed that matrices with entries $\pm 1$ (called Eq.~3 here), or with entries $\sqrt{3}\cdot\{-1, 0, +1\}$ with probabilities $1/6$, $2/3$, and $1/6$ (Eq.~4), work as well as Gaussian ones in that analysis and are cheaper to apply \cite{achlioptas2003}. Li, Hastie, and Church proposed very sparse matrices, whose entries are each nonzero with probability $1/\sqrt{d}$ \cite{li2006}, which makes them sparser than Eq.~4 when $d > 9$.

The assignment applied Eq.~3 and Eq.~4 to five small data sets at $k = 2$ and trained multiclass support vector machines (SVMs) with four kernels on the projected data. The submitted report found Eq.~4 more accurate on three of five data sets and drew conclusions about kernels and data sets from the table. That table rested on one random matrix per distribution, one train/test split, and kernels chosen on the test set.

This paper reports the assignment's result with its draw-to-draw spread, and a new study that asks whether sparse random projections preserve class structure as well as dense ones. It makes three contributions: the corrected assignment results (Section~\ref{sec:assignment}); a paired study of matrix sparsity and target size, with tuned kernels and principal component analysis (PCA) and all-feature baselines (Sections~\ref{sec:design} and~\ref{sec:classification}); and a measurement of distance distortion against its exact second moment and per-pair bounds (Section~\ref{sec:distortion}).

\emph{Revision note.} The original coursework report (spring 2026) compared Eq.~3 and Eq.~4 on one draw. This revised version (fall 2026) replaces that comparison with 100 paired draws, corrects the explanation for the Iris Eq.~4 result, and fixes a few table entries, while also adding the study of sparsity, target size, tuned kernels, and distortion.
